\subsection{Smith, Isayev, Roitberg (2017): ANI-1 transferable NNP for organics}
\label{sec:appendix-example-smith2017}

\paragraph{Q1 -- Bibliographic identification.}
Smith, Isayev, and Roitberg introduce the ANI (ANAKIN-ME)
class of high-dimensional neural-network potentials, in which
atom-type-specific feed-forward NNs take a modified
Behler--Parrinello atomic environment vector (AEV) as input.
The flagship ANI-1 potential is a single transferable model
covering organic molecules made of H, C, N, and O. Published in
\emph{Chem.\ Sci.} 8, 3192--3203 (2017).
\lstinputlisting[language=turtle]{artifacts/kg/listings/smith2017/q01-bibliographic.ttl}

\paragraph{Q2 -- Material system.}
The chemical scope is closed-shell, neutral, singlet organic
molecules over the four-element space $\{$H, C, N, O$\}$, with
training molecules drawn from the GDB-11 enumeration of up to
8 heavy atoms. The benchmark molecules range from 10 to 54
atoms, including pharmaceutical compounds (Retinol, Fentanyl,
Lisdexamfetamine). The protocol's ``one $\MaterialSystemC$ per
paper'' assumption is a tight fit here: we encode this as a
single molecular material with the chemical scope captured in
$\dlRole{materialClass}$ rather than enumerating every benchmark
molecule. No Wikidata entity is used, since the system is a
chemical space rather than a single substance.
\lstinputlisting[language=turtle]{artifacts/kg/listings/smith2017/q02-material.ttl}

\paragraph{Q3 -- Reference calculation method.}
Reference data are produced from molecular DFT calculations.
\lstinputlisting[language=turtle]{artifacts/kg/listings/smith2017/q03-reference-method-type.ttl}

\paragraph{Q4 -- Reference settings.}
DFT calculations are performed with Gaussian~09 using the
$\omega$B97X range-separated hybrid meta-GGA functional and the
6-31G(d) Pople basis set on an ultra-fine integration grid.
There is no Brillouin-zone sampling (molecular calculations);
all-electron Gaussian-basis methodology obviates the explicit
plane-wave $\dlRole{energyCutoff}$ slot. Frozen-core treatment
is not explicitly reported.
\lstinputlisting[language=turtle]{artifacts/kg/listings/smith2017/q04-reference-settings.ttl}

\paragraph{Q5 -- Training dataset.}
The ANI-1 dataset comprises about 17.2~million molecular
conformations generated from $\sim 57\,951$ small molecules
(GDB-11, up to 8 heavy atoms, fluorine removed). 80\% of the
data points are used for training, 10\% for validation, and
10\% for the held-out test set. Energies (single-point) are the
sole covered property; atomic forces are not part of the
training labels.
\lstinputlisting[language=turtle]{artifacts/kg/listings/smith2017/q05-dataset.ttl}

\paragraph{Q6 -- Sampling strategies.}
Two strategies are combined: configurational sampling from the
GDB-11 chemical database (SMILES converted to 3D structures
with RDKit, then DFT-optimised), and conformational
\emph{Normal Mode Sampling} (NMS), in which the equilibrium
structure is displaced along normal-mode coordinates with
random magnitudes drawn from a harmonic-temperature
distribution; $K = S(3N-6)$ structures are generated per
molecule, with $S$ tuned per GDB subset.
\lstinputlisting[language=turtle]{artifacts/kg/listings/smith2017/q06-sampling.ttl}

\paragraph{Q7 -- MLIP method, components, implementation.}
The method is ANI: an HDNNP with atom-type-specific subnets
taking modified Behler--Parrinello AEVs as input (radial
shifted-Gaussian channels, angular shifted-cosine channels with
an additional radial-shell factor; channels split per atom-type
pair). The training loss is the exponential cost
$C = \tau \exp\bigl(\tau^{-1} \sum_j (E^{\mathrm{ANI}}_j -
E^{\mathrm{DFT}}_j)^2\bigr)$ with $\tau = 0.5$. The training
algorithm is ADAM with mini-batches of 1024~molecules and
max-norm regularisation. The implementation is the in-house
GPU-accelerated NeuroChem package (with the AEVLib library for
descriptor computation).
\lstinputlisting[language=turtle]{artifacts/kg/listings/smith2017/q07-method.ttl}

\paragraph{Q8 -- Hyperparameter settings.}
Reported hyperparameter settings: radial cutoff
$R_C^{\mathrm{rad}} = 4.6$~\AA{}, angular cutoff
$R_C^{\mathrm{ang}} = 3.1$~\AA, AEV with 32 radial shifts and
$8\times 8$ angular shifts per element pair (768 total elements
for 4 atom types), per-element pyramidal NN of architecture
$768{:}128{:}128{:}64{:}1$ ($\sim 124\,033$ parameters per
element) with Gaussian hidden activations and a linear output;
ADAM with initial learning rate $10^{-3}$.
\lstinputlisting[language=turtle]{artifacts/kg/listings/smith2017/q08-hyperparameter-settings.ttl}

\paragraph{Q9 -- MLIP run and trained model.}
A single training run produces ANI-1, the flagship potential of
the paper. The optimisation iterates 6 lr-decay restarts on top
of ADAM. No active-learning, ensembling, or per-element
fold-ensemble structure is reported.
\lstinputlisting[language=turtle]{artifacts/kg/listings/smith2017/q09-run-and-model.ttl}

\paragraph{Q10 -- Benchmark results.}
Headline accuracy: total-energy RMSEs (vs.\ $\omega$B97X/6-31G(d))
of 1.2/1.3/1.3~kcal/mol on the training/validation/held-out
test splits, and 1.9~kcal/mol on the GDB-10 transferability
test (134 molecules with 10 heavy atoms, outside the training
distribution; the per-conformation correlation plot reports
1.8~kcal/mol on 8245 conformations). For relative energies
within 30~kcal/mol of the minimum, ANI-1 reaches 0.6~kcal/mol
RMSE; semi-empirical baselines DFTB/PM6/AM1 are at 2.4/3.6/4.2.
\lstinputlisting[language=turtle]{artifacts/kg/listings/smith2017/q10-benchmark-results.ttl}

\paragraph{Q11 -- Computational resources.}
\emph{Not reported quantitatively.} The paper states
qualitatively that NeuroChem accelerates training, testing, and
inference on GPUs via CUBLAS, but it gives no GPU hours,
training duration, hardware spec (model number), peak memory,
or per-atom inference time.
\lstinputlisting[language=turtle]{artifacts/kg/listings/smith2017/q11-resources.ttl}

\paragraph{Q12 -- Gaps.}
Beyond Q11 (no quantitative compute metadata), the following
ontology-expressible items are not reported: an explicit
$\dlRole{frozenCore}$ treatment for the all-electron
6-31G(d) reference; the NeuroChem package version; explicit
recording of forces among $\dlRole{coversProperty}$ (the paper
trains on energies only, so this is a true absence rather than
an omission); per-element data-set sizes within the 17.2~M
total. The ``one material per paper'' assumption is also
stretched here: ANI-1 covers a chemical space, not a single
material, so $\dlRole{sameAsWikidata}$ is omitted intentionally.
